\section{Wild reducible residual inertia at $p=5$}\label{sec:wild-reducible}
We now treat nontrivial residual wild inertia. The high line is then intrinsic: it is the unique fixed line of the residual wild group. Throughout this section $e=20$, $N=24$, $s=120$, and $v_L(5)=480$. A common finite twist normalizes a nonsquare finite determinant so that $c_i\equiv1\pmod4$. A square determinant has already been treated in Proposition~\ref{prop:square-determinant}.

\subsection{Parameters and the two possible matrix groups}
By Lemma~\ref{lem:prime-residual-comparison}, different special characters force $B_i=5b_i$. Equality of the high-line characters and finite determinants gives
\[
 b_i\in\{1,3,5,7,9,11,13,15,17,19\},\qquad
 b_1\equiv b_2\pmod4.
\]
Indeed $c_i-B_i$ is divisible by six, so $b_i$ is odd, and equality of parameters would recover $c_i$. Both parameters cannot be below 20, since the only possibilities are 5 and 15. The boundary 20 is impossible.

If the parameters straddle 20, their high quotients have rational first-level torsion: in the smaller-parameter quotient the old image is a stable low line distinct from its high line. A $5$-group acting on both lines is trivial. Lemma~\ref{lem:one-prime-parameter}, applied to the quotient $25$-comparison, forces the larger $b$ to satisfy $5\mid24-b$. The congruence modulo four leaves exactly
\begin{equation}\label{eq:crossing-pairs}
                         (B_1,B_2)=(5,45),(15,95).
\end{equation}
Their quotient parameters are 25 and 75, with exchanged high lines. If both $B_i>20$, their high quotient parameters are $5(24-b_i)$. Their old first-level images are the common quotient high line. A rational quotient first level would force both $24-b_i$ divisible by five by Lemma~\ref{lem:one-prime-parameter}; the possibilities 9 and 19 for $b_i$ are incongruent modulo four. Thus both original and quotient residual actions are wild.

In the both-large case choose an integral tame eigenbasis with high line $\F_5e_1$. Write the wild image at level 125 as $W$ and its congruence layers as $W_1,W_2$. The functions
\[
                 v_0=\rho_{12}\bmod5,
          \qquad w_0=\rho_{21}/5\bmod5
\]
are additive and nonzero; the second is the quotient residual character. Their tame weights are opposite primitive fourth roots. Thus $w_0$ cannot factor through $v_0$, and the first layer contains $E_{21}$ by a tame projector. Conjugating by $1+tE_{12}$ for $t=1,-1$ gives $E_{12}$ and $\operatorname{diag}(1,-1)$. Fifth powers give $W_1\subset W_2$. Consequently
\[
                  \mathfrak{sl}_2(\F_5)\subset W_1\subset W_2.
\]
The possibility $W_1=W_2=\mathfrak{sl}_2$ is excluded as follows. On the first congruence subgroup the determinant modulo 125 factors through an invariant linear functional on $\mathfrak{sl}_2$. Tame conjugation kills its two off-diagonal values; conjugating $E_{21}$ by $1+E_{12}$ kills the diagonal value. Hence the determinant is trivial there. Its total image has order at most five, so $\mu_{25}\subset K$ and its remaining cyclotomic field equals the residual field $R/L$. The Cohen-field calculation gives break 576. But a low first-level coordinate gives $b(R/L)\le600-6B_i<576$, a contradiction.

There remain two actual groups. Put
\[
 \mathcal I=\left\{
 \begin{pmatrix}1+5a&v\\5w&1+5d\end{pmatrix}:
 a,d,w\in\Z/25,\ v\in\Z/125\right\},
\]
and write each entry in base-five digits, for example $a=a_0+5a_1$ and $v=v_0+5v_1+25v_2$. Then
\begin{equation}\label{eq:two-iwahori-groups}
 W=\mathcal I\quad\hbox{or}\quad
 W=\{\rho\in\mathcal I:a_0+d_0=v_0w_0\}.
\end{equation}
For the first case this follows from the order $5^{1+4+4}=5^9$. In the second, the determinant modulo 25 is trivial on the first layer; its remaining character factors through the residual $v_0$, whose nontrivial tame weight excludes an invariant determinant. Thus the displayed determinant relation holds, and the orders $5^{1+3+4}=5^8$ agree. Notice that both groups retain independent coordinates $v_0,a_0,w_1$ over the quotient residual field, and independent coordinates $w_0,d_0,v_1$ over the original residual field. This fact allows a single point argument to treat both groups.

\subsection{Two fixed-coefficient lemmas}
\begin{lemma}\label{lem:common-high-numbers}
Suppose an elementary abelian degree 25 extension has two coordinate rows
\[
 A_i=5m_i,\qquad D_{0,i}=125m_i,\qquad
 D_{1,i}=15000-625m_i,\qquad v(5)=60000,
\]
with a common high-displacement line, where $1\le m_1<m_2\le19$ and both $m_i$ are prime to five. Then that line is last, and the two lower breaks and the larger upper break are
\[
 t=24m_1,\qquad u=15000-630m_2-96m_1,
 \qquad C=3000-126m_2.
\]
\end{lemma}
\begin{proof}
A single break $h$ is impossible. If the first coordinate change were below $D_{0,i}$, a nonzero additive polynomial of degree at most five would vanish on all 25 injective leading coefficients. At $D_{0,i}$ a unique monomial could not vanish precisely on the high line. A tie is therefore necessary, forcing $h=(D_{0,i}-A_i)/5=24m_i$ for both models. These values differ.

For two breaks $t<u$, a last line different from the high line would give larger upper break $24m_i$ in both models by Lemma~\ref{lem:opposite-line-conductor}. Hence the high line is last. Put $r_i=24m_i$. The expansion~\eqref{eq:first-change}, with first ordinary exponent $J_i$, gives
\[
\begin{array}{c|c|c}
 &J_i&\text{last-break equality}\\\hline
 t<r_i&5m_i+4t&u+4t=15000-630m_i\\
 t>r_i&125m_i-t&u-t=15000-750m_i.
\end{array}
\]
In the first row the two earlier terms must cancel; in the second the ordinary term is uniquely first. At the later break it is uniquely first in both cases. The same strict row in both models forces equal parameters. The mixed strict case gives $t=150m_1-126m_2<24m_1$, contrary to its assumed position. At $t=r_2$, the first model gives $u=15000-750m_1+24m_2$. For the second model, $A_2+5u-D_{1,2}=750(80-5m_1+m_2)>0$, because $m_1<m_2\le19$. Thus the ordinary term controls the later displacement. The boundary requires $J_2\ge101m_2$, but that later displacement forces $J_2=750m_1-649m_2$, implying $m_1\ge m_2$. Thus $t=r_1$, and the larger parameter supplies the asserted formulas.
\end{proof}

\begin{lemma}\label{lem:large-gap}
Let $T/S$ be totally ramified elementary abelian of degree 25, with a line $H$ in its group. Suppose coordinates have values $A_i>0$, $5\mid A_i$, and displacements $D_{1,i}$ on $H\setminus\{0\}$ and $D_{0,i}$ off $H$, with
\[
                    A_i<D_{0,i}<D_{1,i}<v_T(5).
\]
Put $\Delta_i=D_{0,i}-A_i$. If the cyclic quotient killing $H$ has conductor $z$ satisfying
\[
 z>\frac9{25}\min_i\Delta_i,\qquad
 z>\frac15\max_i\Delta_i,
\]
then every gap $D_{1,i}-D_{0,i}$ is the same.
\end{lemma}
\begin{proof}
In the expansion in a uniformizer with fixed coefficient representatives, every five-divisible exponent has its first possible change at degree at least $A_i+5h$ at break $h$. The first nonzero ordinary exponent $J_i$ changes at $J_i+h$ with nonzero coefficient. This holds also when $25\mid A_i$; terms with five-divisible binomial coefficients are beyond the stated mixed-characteristic cutoff.

If there is a single break $z$, the second inequality places all divisible terms after $D_{0,i}$. An ordinary term must give $J_i+z=D_{0,i}$, and it is nonzero for every nonidentity element, including those of $H$. This contradicts $D_{1,i}>D_{0,i}$.

Suppose the lower breaks are $t<u$, with last line $V\ne H$. Choose $i$ minimizing $\Delta_i$. The displacement on $V$ gives $u\le\Delta_i$. If $t>\Delta_i/5$, the ordinary term forced at $u$ would already give $J_i+t<D_{0,i}<D_{1,i}$ on $H$, a contradiction. Thus $t\le\Delta_i/5$, and the quotient conductor is $(u+4t)/5\le9\Delta_i/25$, again impossible. Hence $V=H$ and $t=z$. At $t$, every divisible term is later than $D_{0,i}$, so $J_i+t=D_{0,i}$. At $u$ these terms increase at least five times as fast as the ordinary term; therefore $D_{1,i}=J_i+u$. The common gap is $u-t$.
\end{proof}

\subsection{A selected point excludes both large parameters}
\begin{proposition}\label{prop:both-large-excluded}
A different-special-character comparison cannot have $B_1,B_2>20$.
\end{proposition}
\begin{proof}
The possible different-character parameter pairs, after ordering, are
\[
 (5,9),(5,13),(5,17),(7,15),(7,19),
 (9,13),(9,17),(11,15),(11,19),(15,19).
\]
Here the entries are $b_i=B_i/5$. This list follows by combining $b_1\equiv b_2\pmod4$, $c_i\equiv1\pmod4$, $c_i\equiv5b_i\pmod6$, and equality of $c_i-5b_i$ modulo 24, then removing the two pairs $(7,11),(13,17)$ giving equal unordered orbits $\{c_i,5c_i\}$. The involution $(b_1,b_2)\mapsto(24-b_2,24-b_1)$ pairs the ten rows. We may choose the original or high-quotient orientation to obtain
\begin{equation}\label{eq:large-five-pairs}
                    (5,9),(5,13),(5,17),(7,15),(11,15).
\end{equation}
This change of orientation will use a selected quotient point only, not an extra level of quotient torsion.

For the original orientation set $P^+=e_1/125$, $P=5P^+$, $E_h=L(P)$, $S=R_oE_h$, and $T=R_oL(P^+)$, where $R_o$ is the field of $v_0$. In either group~\eqref{eq:two-iwahori-groups}, their stabilizers respectively impose
\[
 a_0=w_0=0;\qquad v_0=a_0=w_0=0;\qquad v_0=a=w=0.
\]
Their degrees over $L$ are 25,125,3125. On the subgroup for $S$, $(a_1,w_1)$ are independent additive coordinates and their joint kernel fixes $T$. Thus $T/S$ is elementary abelian of degree 25, with high line $w_1=0$.

In the quotient orientation use the actual image $M[125]/H$, of type $\Z/125\oplus\Z/25$, and choose $P^+=q(e_2/125)$. Its order is 125 and its first-level image is the quotient high point. In the reordered quotient lattice $(e_2,e_1/5)$ the matrix is
\[
                 \begin{pmatrix}1+5d&w\\5v&1+5a\end{pmatrix}.
\]
Its first column is determined modulo 125 by the original module. The stabilizers of $P$, $S$, and $T$ impose respectively
\[
 v_0=d_0=0;\qquad w_0=v_0=d_0=0;\qquad
 w_0=v_0=v_1=d_0=d_1=0.
\]
The same degrees result in both groups, and $T/S$ has additive coordinates $(d_1,v_1)$ and high line $v_1=0$. The trace relation in~\eqref{eq:two-iwahori-groups} only determines the unused diagonal digit in each of these calculations.

For either orientation the two successive high inverse polygons are single segments. A high first-level value $5(24-b_i)$ becomes $(24-b_i)/5$ and then $(24-b_i)/125$. Hence in $T$ the coordinates have
\begin{equation}\label{eq:large-point-rows}
 A_i=25(24-b_i),\quad D_{0,i}=3125b_i,\quad
 D_{1,i}=15625(24-b_i),\quad v_T(5)=1500000.
\end{equation}
These displacements are the translating first-level values multiplied by 3125.

We compute a lower bound for the quotient conductor killing the high line. In the opposite orientation take a primitive low 25 point. Its field $Z$ has degree 125 over $L$ and degree 25 over its residual field $R_*$. Its two rows are those of Lemma~\ref{lem:common-high-numbers}, with parameter pairs
\[
                    (15,19),(11,19),(7,19),(9,17),(9,13).
\]
The diagonal and off-diagonal tame weights differ, so there are two ramification breaks. For pairs prime to five, the common high line is last by Lemma~\ref{lem:opposite-line-conductor}. For $(15,19)$ this conclusion still holds, as follows. If the high line were not last, the prime row would give $u+4t=120\cdot19$. In the other row $A=75$, and the first possible changes have degrees
\[
                \min\{75+25h,N_1+5h,J+h\},
\]
where $N_1>75$ is the first exponent of five-adic order one and $J>75$ the first ordinary exponent. At $t$ at least two terms must tie, because the required high displacement is larger than the low displacement at $u$. A tie of the last two terms gives $J=N_1+4t$ and low displacement $N_1+120\cdot19>1875$. A tie of the first and last gives $J=75+24t$ and an even larger displacement. A tie of the first two with the last later gives $N_1=75+20t$, while at $u$ both divisible terms and the ordinary term are later than 1875. All cases contradict the prescribed low displacement. Thus the high line is last.

Let $t'$ and $C'$ be the diagonal and off-diagonal conductors of $Z/R_*$. The resulting values are
\[
\begin{array}{c|c|r}
 (m_1,m_2)&t'&C'\\\hline
 (15,19)&24\le t'\le336,\ 24\mid t'&606\\
 (11,19)&264&606\\
 (7,19)&168&606\\
 (9,17)&216&858\\
 (9,13)&216&1362.
\end{array}
\]
The four prime rows follow from Lemma~\ref{lem:common-high-numbers}. For the mixed row, if $t'>24\cdot19$, the first ordinary exponents in both rows force incompatible values $u-t'=15000-750m_i$; equality at $24\cdot19$ gives the same boundary contradiction as in that lemma. Thus $t'<456$, and the prime row gives $u=3030-4t'$. In the row $m_1=15$, one has $5u+75>D_1$ because $t'<456$, and its forced ordinary exponent is $J=5625-u$. Its change at $t'$ is later than 1875, since $J+t'-1875=720+5t'>0$. If $t'\ge360$, every divisible term is also later than 1875: the leading degree 75 has order two, and all order-one degrees exceed 75. Therefore $t'<360$. The diagonal tame weight gives $24\mid t'$, and the upper conversion gives $C'=606$.

In the original orientation, on the subgroup over $R_*$ the three coordinates $(v_0,a_0,w_1)$ are independent additive characters in both groups~\eqref{eq:two-iwahori-groups}. The field $Z$ has characters $(a_0,w_1)$, and $S$ has $(v_0,a_0)$. If $\sigma$ is the conductor of $v_0$, conjugation by an opposite residual unipotent sends $a_0$ to $a_0\pm t v_0$. Invariance of conductors gives $\sigma\le t'$, and the distinct tame weights exclude equality through the injective leading-coefficient map. Thus $0<\sigma<t'<C'$. The actual compositum $SZ/R_*$ has these three independent characters, so the conductor of $w_1$ restricted to $S$ is
\[
 z=\psi_{S/R_*}(C')
   =\sigma+5(t'-\sigma)+25(C'-t')>25C'-24t'.
\]
It is the quotient of $T/S$ killing its high line. For the quotient orientation the same calculation uses $(w_0,d_0,v_1)$, with $Z$ defined by $(d_0,v_1)$ and $S$ by $(w_0,d_0)$. Thus it proves the same bound without any hypothetical quotient 125 comparison.

Substitution gives
\[
\begin{array}{c|r|r|r}
 (b_1,b_2)&z>&9\min_i(D_{0,i}-A_i)/25&\max_i(D_{0,i}-A_i)/5\\\hline
 (5,9)&7086&5454&5550\\
 (5,13)&8814&5454&8070\\
 (5,17)&11118&5454&10590\\
 (7,15)&16266&7722&9330\\
 (11,15)&28866&12258&9330.
\end{array}
\]
All hypotheses of Lemma~\ref{lem:large-gap} hold. It forces equality of the gaps in~\eqref{eq:large-point-rows}, but these are $375000-18750b_i$, different for $b_1<b_2$. This excludes both actual groups and every both-large pair.
\end{proof}

\subsection{The two crossing parameters}
\begin{lemma}\label{lem:opposite-twenty-five}
Suppose a comparison of $25$-torsion has rational first-level modules over $L$, parameters 25 and 75, and exchanged high lines $\F_5e_1,\F_5e_2$. Its image is $1+5M_2(\F_5)$, and its coordinate characters have breaks
\[
             c(A_{21})=6,\quad c(A_{22})=24,
             \quad c(A_{12})=66,\quad c(A_{11})=72.
\]
\end{lemma}
\begin{proof}
The degree-five point-field exclusions in the proof of Theorem~\ref{thm:tame-five-recovery} use only the $25$-comparison: the coordinate values are $(19/5,3/5)$ and $(1/5,9/5)$, and the six pairs of exact degree-five breaks displayed there all disagree. Thus both point fields $E_i$ have degree 25. Tame projection puts both off-diagonal axes in the image and gives nonzero projections to both diagonal coordinates. Its dimension is at least three.

The four coordinate rows are the first four rows in that theorem's table. In $E_1$, an off-diagonal last line would give last upper break 6 by Lemma~\ref{lem:opposite-line-conductor}, impossible since the diagonal cyclic break is a positive multiple of 24. Thus the diagonal is last, with upper break 72. In $E_2$, a diagonal last line would have break 24; an off-diagonal last line would have break 66. If the matrix image had dimension three, its two diagonal projections would be the same nonzero cyclic character up to scalar, and hence have the same conductor. Neither alternative in $E_2$ can share the diagonal break 72 of $E_1$. Therefore the image is all $M_2(\F_5)$.

We determine the remaining breaks without a computation of any higher torsion. In $E_1$, the smaller lower break satisfies $t\equiv6\pmod{24}$ and $0<t<72$, and the larger is $u=360-4t$. For $t=30$, the second row forces the ordinary exponent $15+4t=135$, divisible by five. For $t=54$, the first row forces the exponent $125-t=71<95$. Both are impossible, so $(t,u)=(6,336)$.

In $E_2$, a diagonal last line would give $(t,u)=(18,48)$. The fourth coordinate row forces the first ordinary exponent 117, which would change at 165 at the later break instead of 1125. If the off-diagonal is last, its upper break is 66 and the diagonal smaller break is 24 or 48. The latter forces the third row's ordinary exponent 77 and a later change 215 instead of 2375. Hence $(t,u)=(24,234)$. Upper conversion gives precisely the four stated cyclic conductors.
\end{proof}

In either crossing~\eqref{eq:crossing-pairs}, the high quotients satisfy this lemma. Their lattice $(e_1/5,e_2)$ has matrix
\[
 \begin{pmatrix}\rho_{11}&5\rho_{12}\\\rho_{21}/5&\rho_{22}\end{pmatrix}.
\]
Rationality of its first level forces $25\mid\rho_{21}$. Its full 25 image makes the four functions $v_0=\rho_{12}\bmod5$, $a_0=(\rho_{11}-1)/5\bmod5$, $d_0=(\rho_{22}-1)/5\bmod5$, and $c=\rho_{21}/25\bmod5$ jointly surjective. The original first layer is upper triangular, with the full diagonal plane by tame projection and $E_{12}$ from fifth powers of residual unipotents. The terminal layer contains this plane and $E_{12}$. If its lower coordinate vanished, the nonzero character $c$ would factor successively through the first layer and residual upper character, both of different tame weights. This is impossible. Hence the terminal layer is all $M_2(\F_5)$, and the full image is
\begin{equation}\label{eq:crossing-matrix}
 W=\left\{\begin{pmatrix}1+5a&v\\25c&1+5d\end{pmatrix}:
 a,d\in\Z/25,\ v\in\Z/125,\ c\in\F_5\right\},
 \qquad |W|=5^8.
\end{equation}
The order proves equality with the displayed containing group.

Let $R$ be the original residual field and $D$ the cyclic field of $d_0$. In the first crossing the original residual upper entry is the quotient character $A_{12}$ and $d_0$ is $A_{22}$; their breaks are 66 and 24. In the second crossing the parameter 25 quotient is the other model, so reordering its high basis identifies these with $A_{21}$ and $A_{11}$; their breaks are 6 and 72.

A low 25 point has field $E_\ell/L$ of degree 125, with $E_\ell/R$ elementary abelian of degree 25 and translation coordinates $(v_1,d_0)$. Its rows are those of Lemma~\ref{lem:common-high-numbers}, with $(m_1,m_2)=(1,9)$ or $(3,19)$. The quotient killing the high line is the actual field $DR/R$. In the second crossing that lemma gives its break 72, whereas the actual distinct fields $R,D$ over $L$ have breaks 6 and 72, giving restricted break $6+5(72-6)=336$. This contradiction excludes $(15,95)$.

\subsection{The terminal determinant excludes the last crossing}
Theorem~\ref{thm:cyclotomic-exact} gives the last upper break of $K(\mu_{125})/K$ as $44-d\ge20$, where $d=v_K(\mathfrak D_{K/K_0(\mu_5)})\le24$. Its last break after the tame extension $L/K$ is therefore at least $480$.

Assume now the remaining crossing $(5,45)$. Let $F$ and $F_2$ be the original 125 and 25 division fields. In~\eqref{eq:crossing-matrix}, the four cyclic fields defined by $(v_0,a_0,d_0,c)$ have breaks $(66,72,24,6)$ over $L$, by Lemma~\ref{lem:opposite-twenty-five}. Put $S=R A$, where $A$ is the field of $a_0$. The common-high lemma gives lower breaks 24 and 9234 for $E_\ell/R$, hence upper breaks 24 and 1866 for its characters $(d_0,v_1)$. The character $a_0$ restricted to $R$ has conductor $66+5(72-66)=96$. Thus $F_2/R$ has upper breaks 24,96,1866 and lower breaks 24,384,44634.

For precision we also compute the passage to $L$. The nonidentity elements of $\Gal(F_2/R)$ have these three lower breaks with multiplicities 100,20,4. If $M$ is the maximal lower break in a nontrivial residual coset, the quotient-average formula gives
\[
 \frac{100\min(24,M)+20\min(384,M)+4\min(44634,M)+M}{125}=66.
\]
Its unique solution is 234. Indeed the left side is strictly increasing and equals $24+(M-24)/5$ between 24 and 384. The other elements of the coset have the minimum of $M$ and the corresponding subgroup break, by the ultrametric inequality and maximality. Therefore
\begin{equation}\label{eq:crossing-filtrations}
\begin{array}{c|c|c}
 \text{extension}&\text{lower breaks}&\text{upper breaks}\\\hline
 F_2/L&24,234,384,44634&24,66,72,426\\
 F_2/S&24,44634&24,8946\\
 F_2/E_\ell&384&384.
\end{array}
\end{equation}
The last two rows also follow directly from the actual stabilizers: over $S$ one has $v_0=a_0=0$ and free $(d_0,v_1)$; over $E_\ell$ one has $v_0=v_1=d_0=0$ and free $a_0$.

We bound the two diagonal characters on the terminal group $\Gal(F/F_2)=1+25M_2(\F_5)$. Choose a high 25 inverse in the first model's high inverse coset. Its value is 19; the second model's value is 3. Tame projection to the first eigenline preserves that coset, so the point and a $125$-torsion lift may be chosen in the same integral tame eigenline. Both next inverse polygons are single segments, giving values $19/25$ and $3/25$.

For the 125 point $P_h^+$ put $T_h=R L(P_h^+)$. Its stabilizer is $v_0=a=c=0$, so $[T_h:L]=625$ and $T_h/S$ is elementary abelian with independent additive coordinates $(a_1,c)$. These assertions follow by multiplying the matrices on the subgroup $v_0=a_0=0$. The character $c$ has conductor 6 over $S$, since both breaks of $S/L$ are larger. In $T_h$ the second coordinate has value 75, while every nonzero high translation has displacement $625\cdot75=46875$. Hence its lower break is at most 46800. If the high line is last, the first break is 6 and the last upper break is at most $(46800+24)/5$; if it is not last, the $c$ conductor 6 is already the largest. Thus in every case
\[
                         c_S(a_1)\le46824/5.
\]
Restriction to $F_2$, using~\eqref{eq:crossing-filtrations}, gives
\begin{equation}\label{eq:crossing-diagonal-bounds}
 c_{F_2}(a_1)\le44634+25(46824/5-8946)=55104.
\end{equation}
No disjointness assertion is needed for this upper restriction bound.

For a low 125 lift $P_\ell^+$, its field has degree 3125 over $L$ and is elementary abelian of degree 25 over $E_\ell$, with additive coordinates $(v_2,d_1)$. The first-model coordinate has value 5. Every element with $d_1\ne0$ displaces it by 3125, so every such lower break is at most 3120. Averaging over the high line gives $c_{E_\ell}(d_1)\le3120$. Restricting through the cyclic last row of~\eqref{eq:crossing-filtrations} gives
\[
                    c_{F_2}(d_1)\le384+5(3120-384)=14064.
\]
Both characters are the actual terminal diagonal coordinates. The first is already additive on the subgroup for $S$, and the second on the subgroup for $E_\ell$; this verifies their identities before applying any restriction formula.

The actual determinant modulo 25 is $a_0+d_0$, whose conductor over $L$ is 72. Hence $\mu_{25}\not\subset K$. By Theorem~\ref{thm:cyclotomic-exact}, the last cyclotomic break over $L$ is at least 480, greater than the last break 426 of $F_2/L$. Its terminal restriction therefore has conductor at least
\[
                  \psi_{F_2/L}(480)=44634+625(480-426)=78384.
\]
On the terminal group, however, the same determinant is $1+25(a_1+d_1)$, whose conductor is at most 55104 by the two proved bounds. This contradiction excludes the last crossing.

\begin{theorem}\label{thm:wild-five-recovery}
Two potentially supersingular height-two groups over a field with algebraically closed residue field of characteristic five, $e_K=20$, and $\mu_5\subset K$, having isomorphic 125-torsion and nontrivial residual wild inertia, have equal finite special characters.
\end{theorem}
\begin{proof}
A square finite determinant is covered by Proposition~\ref{prop:square-determinant}. Otherwise the parameter reduction leaves precisely the crossings~\eqref{eq:crossing-pairs} or both-large parameters. The preceding arguments exclude the crossings, and Proposition~\ref{prop:both-large-excluded} excludes both-large parameters for both possible congruence images. These cases exhaust the different-character assumption.
\end{proof}
